% Effective Spacetime Geometry from Admissible Non-Injective Projection
% Cosmochrony companion paper Q6b

\documentclass[11pt,a4paper]{article}

\usepackage{amsmath,amssymb,amsthm}
\usepackage{enumerate}
\usepackage{hyperref}
\usepackage{doi}
\usepackage[numbers,sort&compress]{natbib}

\newtheorem{theorem}{Theorem}[section]
\newtheorem{lemma}[theorem]{Lemma}
\newtheorem{proposition}[theorem]{Proposition}
\newtheorem{corollary}[theorem]{Corollary}
\newtheorem{definition}[theorem]{Definition}
\newtheorem{remark}[theorem]{Remark}

\title{Effective Spacetime Geometry\\
from Admissible Non-Injective Projection}

\author{J\'er\^ome Beau\\
\small Independent Researcher\\
\small \texttt{jerome.beau@cosmochrony.org}}

\date{2026}

\begin{document}

  \maketitle

  \begin{abstract}
    Continuing from the gauge structure established in the companion paper Q6a~\cite{Beau2026q6a}, we derive effective
    spacetime geometry from the admissibility framework of Foundation~M~\cite{Beau2026m} by importing the continuum
    limit and geometric-extraction theorems of the companion paper Q5b~\cite{Beau2026q5b}.
    Under the lifting hypothesis~[H-lift] of Q5b and the Mosco hypotheses of Q5a~\cite{Beau2026q5a}, the admissibility
    filter $\Pi_q$ converges to an effective operator $\mathcal{L}_{\mathrm{eff}}$ on $\mathbb{R}_\tau \times
    \mathrm{Heis}_3(\mathbb{R})$ whose principal symbol defines a four-dimensional effective metric tensor
    $g^{\mu\nu}(x) \propto A^{\mu\nu}(x)$ of Lorentzian signature~$(-,+,+,+)$
    (Q5b Theorems~5.2 and~6.1).
    In the presence of a localised stationary obstruction to the admissibility flow with spherical symmetry, flux
    conservation within this effective geometry forces the Schwarzschild metric as the unique stationary exterior
    solution, with the horizon appearing as a degeneracy of the principal symbol rather than as a singularity of the
    underlying admissible structure.
    The Einstein equations emerge as consistency conditions of the emergent geometry, with the Einstein\textendash
    Hilbert action derived rigorously from the spectral entropy functional in the Gravity
    paper~\cite{Beau2026h}.
    The present paper closes the chain $\Pi_q \to \mathcal{L}_{\mathrm{eff}} \to g^{\mu\nu} \to G_{\mu\nu}$.
    No dynamical substrate is required at any step.
  \end{abstract}

  \section{Introduction}
    \label{sec:intro}

    \subsection*{Position in the programme}

      The Cosmochrony programme~\cite{Cosmochrony} derives physical structure from a static relational substrate
      $\chi$ through a non-injective projection $\Pi$ onto an observable space.
      The companion paper Q6a~\cite{Beau2026q6a} derives gauge charges, the gauge group, and Yang\textendash Mills
      dynamics from the fibre structure of~$\Pi$.
      The present paper addresses the complementary macroscopic question: what effective geometry is selected when the
      bounded admissibility constraints of Foundation~M~\cite{Beau2026m} are imposed on a spatially varying
      projection?

      The geometric chain is now in place.
      Q5a~\cite{Beau2026q5a} establishes, under Hypotheses~H1,~H-w,~H-E1 and Conjecture~C, that the discrete admissible
      Dirichlet forms on the Cayley graph $\mathrm{Cay}(G_q, S_q)$ Mosco-converge to a one-dimensional shadow operator
      $\mathcal{L}_\Pi = -A\,\partial_x^2$ on $L^2(\mathbb{R})$.
      Q5b~\cite{Beau2026q5b} promotes this one-dimensional shadow to a four-dimensional effective geometry: the BFS
      shell stratification of $G_q$ converges, in the Gromov\textendash Hausdorff sense, to the
      Carnot\textendash Carath\'eodory sphere foliation of $\mathrm{Heis}_3(\mathbb{R})$ with homogeneous dimension
      $D_{\mathrm{hom}} = 4$, and under the lifting hypothesis~[H-lift] the principal symbol of the effective operator
      extends to a Lorentzian metric tensor on $\mathbb{R}_\tau \times \mathrm{Heis}_3(\mathbb{R})$.

      The task of the present paper is to close the chain from this effective metric to the Einstein equations of
      general relativity.
      Three results are obtained: (i) the Schwarzschild metric as the unique stationary, spherically symmetric
      exterior solution compatible with the admissibility flow (§\ref{sec:schwarzschild}); (ii) a structural
      interpretation of the horizon as a projection degeneracy (§\ref{sec:horizon}); (iii) the identification of the
      Einstein equations as consistency conditions of the emergent geometry, via the spectral entropy functional of the
      Gravity paper~\cite{Beau2026h} (§\ref{sec:einstein}).

    \subsection*{What is and is not derived in this paper}

      This paper derives the \emph{structure} of effective spacetime geometry from the admissibility framework.
      It does not derive the precise numerical value of Newton's constant $G_N$, nor does it prove the hypotheses
      inherited from Q5a and Q5b.
      All results are therefore conditional on~[H-lift] (Q5b Hypothesis~4.1) and on the Q5a Mosco convergence
      hypotheses.
      Open directions are collected in Section~\ref{sec:open}.

    \subsection*{Organisation}

      Section~\ref{sec:continuum} reviews the Q5a\textendash Q5b geometric chain and collects the theorems used as
      inputs.
      Section~\ref{sec:geometry} derives the effective metric, imports its Lorentzian signature, establishes the
      Schwarzschild exterior solution, and interprets the horizon.
      Section~\ref{sec:einstein} closes the chain to the Einstein equations.
      Section~\ref{sec:open} identifies the open directions.

  \section{The Q5a\textendash Q5b Geometric Chain}
    \label{sec:continuum}

    \subsection{From the discrete projection to the 1D shadow operator}
      \label{sec:q5a}

      The discrete admissibility filter $\Pi_q$ of Foundation~M~\cite{Beau2026m} acts on the orbit span
      $\Omega_n^{(c)} \subset \mathbb{C}^q$, selecting admissible directions via the Born\textendash Infeld bound
      $A_n \leq c_\chi/\sqrt{\lambda_n}$ (Axiom~A3 of~\cite{Beau2026m}).
      Q5a~\cite{Beau2026q5a} establishes three convergences, conditionally on Hypotheses~H1,~H-w,~H-E1 and
      Conjecture~C:
      \begin{itemize}
        \item[(T1)] a Hilbert inductive limit $\varinjlim_q \mathcal{C}_q \simeq L^2(\mathbb{R})$, where
        $\mathcal{C}_q \simeq L^2(\mathbb{Z}/q\mathbb{Z})$ is the Weil representation space at finite~$q$;
        \item[(T2)] convergence of rescaled Weil generators to the metaplectic generators of
        $\mathrm{Heis}_3(\mathbb{R})$;
        \item[(T3)] Mosco convergence of the admissibility Dirichlet forms $\mathcal{E}_q$ to a second-order operator
        \begin{equation}
          \label{eq:Lpi-shadow}
          \mathcal{L}_\Pi \;=\; -A\,\partial_x^2
          \quad\text{on}\quad L^2(\mathbb{R}),
          \qquad A > 0.
        \end{equation}
      \end{itemize}

      The operator~\eqref{eq:Lpi-shadow} is one-dimensional: it acts on the central coordinate~$x$ of the Weil block,
      and is the operator-theoretic shadow of the homogeneous four-dimensional structure carried by
      $\mathrm{Heis}_3(\mathbb{R})$.
      The passage from this one-dimensional shadow to a four-dimensional effective geometry is the content of Q5b.

    \subsection{From the 1D shadow to a four-dimensional Carnot geometry}
      \label{sec:q5b}

      Q5b~\cite{Beau2026q5b} establishes three further results that together promote~$\mathcal{L}_\Pi$ to a
      four-dimensional effective metric.

      \paragraph{Carnot convergence.}
        By the Bass\textendash Guivarc'h theorem, $\mathrm{Heis}_3(\mathbb{R})$ is a Carnot group of homogeneous dimension
        $D_{\mathrm{hom}} = 4$ (Q5b Theorem~2.1).
        The BFS shell stratification $\{S_n\}$ of $G_q = \mathrm{Heis}_3(\mathbb{Z}/q\mathbb{Z})$ converges, in the
        pre-saturation regime $n \ll q$, to the Carnot\textendash Carath\'eodory sphere foliation of the continuous group
        as a structural consequence of Pansu's asymptotic cone theorem (Q5b Theorem~3.2~\cite{Beau2026q5b}).
        The limiting geometry therefore carries spectral dimension~$4$.

      \paragraph{Lifting hypothesis.}
        Q5b introduces Hypothesis~[H-lift] (Definition~4.1 of~\cite{Beau2026q5b}): the one-dimensional shadow
        $\mathcal{L}_\Pi$ of~\eqref{eq:Lpi-shadow} is the image, under the Schr\"odinger representation~$d\pi_1$, of the
        kinetic sector of the sub-Riemannian Laplacian $\Delta_H = \widetilde{X}^2 + \widetilde{Y}^2$ on
        $\mathrm{Heis}_3(\mathbb{R})$, restricted to the admissible sector.
        This hypothesis is open (Q5b-O1) and is expected to follow from the Plancherel decomposition of
        $L^2(\mathrm{Heis}_3(\mathbb{R}))$; it is the only hypothesis introduced by Q5b beyond those of Q5a.

      \paragraph{Metric tensor from the principal symbol.}
        Under~[H-lift], the full effective operator $\mathcal{L}_{\mathrm{eff}}$ on
        $\mathbb{R}_\tau \times \mathrm{Heis}_3(\mathbb{R})$ has a principal symbol
      \begin{equation}
        \label{eq:principal-symbol}
        \sigma_2(\mathcal{L}_{\mathrm{eff}})(x, k) \;=\; A^{\mu\nu}(x)\,k_\mu k_\nu,
        \end{equation}
        where $A^{\mu\nu}(x)$ is a symmetric tensor with entries $A_\tau$ in the temporal sector and $A_H$ in the
        sub-Riemannian horizontal sector.
        Q5b Theorem~5.2 extracts from this symbol the effective metric
      \begin{equation}
        \label{eq:metric-from-symbol}
        g^{\mu\nu}(x) \;\propto\; A^{\mu\nu}(x),
        \end{equation}
        defined on the non-degenerate part of the sub-Riemannian structure.
        The metric is not an independent dynamical variable: it is a derived object encoding how admissible modes
        propagate in the continuum limit.

      \paragraph{Lorentzian signature.}
        In homogeneous and quasi-isotropic admissible regimes, the bounded Born\textendash Infeld admissibility
        constraint (A3\textendash A4 of~\cite{Beau2026m}) forces $g^{\mu\nu}$ to have Lorentzian signature $(-,+,+,+)$
        (Q5b Theorem~6.1, importing the argument of~\cite{Beau2026c}~Sections~5\textendash 6).
        The signature is selected by exclusion of the three incompatible classes: Riemannian signature produces an
        elliptic principal symbol and contradicts the finite-capacity bound of~A4; ultra-hyperbolic signature fails to
        admit a globally well-posed initial value problem; degenerate signature is excluded at the level of the effective
        metric after incorporation of sub-principal corrections.

  \subsection{Inherited status discipline}
    \label{sec:status}

    The results of the present paper are conditional on the following hypotheses, all inherited from
    Q5a~\cite{Beau2026q5a} and Q5b~\cite{Beau2026q5b} and carried forward without modification:
    \begin{itemize}
      \item H1, H-w, H-E1 (Q5a): embedding, weight stabilisation, and energy control hypotheses governing the Mosco
      limit of~\eqref{eq:Lpi-shadow}.
      \item Conjecture~C (Q5a): discrete Rellich compactness, numerically confirmed at the pair level for
      $q \leq 401$ via the O-series pipeline~\cite{Beau2026a16,Beau2026a28}.
      \item [H-lift] (Q5b~Definition~4.1): identification of $\mathcal{L}_\Pi$ as the kinetic sector of
      $d\pi_1(\Delta_H)$.
    \end{itemize}
    Closing these hypotheses is the subject of open direction~Q5b-O1 (Section~\ref{sec:open}) and would remove all
    conditionality from the results below.

  \section{Effective Spacetime Geometry}
  \label{sec:geometry}

  \subsection{The effective metric}
    \label{sec:metric}

    Under the hypotheses of Section~\ref{sec:status}, Q5b Theorem~5.2 gives the effective metric
    $g^{\mu\nu}(x) \propto A^{\mu\nu}(x)$ on the non-degenerate part of the sub-Riemannian structure, with Lorentzian
    signature~$(-,+,+,+)$ from Q5b Theorem~6.1.
    The metric is defined on the effective continuum $\mathbb{R}_\tau \times \mathrm{Heis}_3(\mathbb{R})$ and
    governs the propagation of admissible modes in the continuum limit of the projection.

    \begin{remark}[Nature of the effective four-dimensionality]
      \label{rem:4d}
      The effective four-dimensionality is not that of a non-degenerate pseudo-Riemannian manifold at the
      principal-symbol level.
      It arises from the homogeneous dimension $D_{\mathrm{hom}} = 4$ of the underlying Carnot structure (Q5b
      Theorem~3.2), which governs both volume growth and spectral scaling.
      The Lorentzian metric extracted from the principal symbol is an effective object, valid after projection onto
      admissible continuum configurations; it is not a smooth metric at every point of the sub-Riemannian space
      $\mathbb{R}_\tau \times \mathrm{Heis}_3(\mathbb{R})$ without further hypotheses.
      Extending the construction to a full $4\times 4$ non-degenerate Lorentzian metric that incorporates the
      central direction~$Z$ is the content of Q5b-O2.
    \end{remark}

  \subsection{Schwarzschild geometry from flux conservation}
    \label{sec:schwarzschild}

    The Lorentzian structure of Section~\ref{sec:metric} is defined in homogeneous and quasi-isotropic admissible
    regimes.
    A localised stationary obstruction to the admissibility flow breaks homogeneity and selects, via flux
    conservation, a specific stationary exterior metric.

    \begin{theorem}[Schwarzschild from bounded admissibility, after~\cite{Beau2026c}]
      \label{thm:schwarzschild}
      Let the effective geometry of Section~\ref{sec:metric} be placed in the presence of a localised stationary
      obstruction to the admissibility flow, with spherical symmetry and no additional microscopic scales.
      Then the unique exterior stationary solution of the effective geometry induced by $\mathcal{L}_{\mathrm{eff}}$
      is the Schwarzschild metric
      \begin{equation}
        \label{eq:schwarzschild}
        ds^2 \;=\; -\!\left(1 - \frac{r_s}{r}\right)dt^2
        + \left(1 - \frac{r_s}{r}\right)^{-1}\!dr^2
        + r^2\,d\Omega^2,
      \end{equation}
      where $r_s$ is an integration constant characterising the strength of the obstruction.
    \end{theorem}

    \begin{proof}
      In the exterior region, the homogeneous equation for the effective operator of Section~\ref{sec:q5b} under
      spherical symmetry and stationarity reduces to the flux conservation equation
      \begin{equation}
        \label{eq:flux-conservation}
        \frac{1}{r^2}\,\frac{d}{dr}\!\left(r^2 A^r(r)\,\frac{d\Phi}{dr}\right) \;=\; 0,
      \end{equation}
      expressing conservation of admissible flux through spheres of radius~$r$.
      The general solution is $\Phi(r) = \Phi_0 - C/r$, giving a $1/r$ decay profile.
      Via the operator-metric correspondence~\eqref{eq:metric-from-symbol}, this profile translates into the metric
      coefficients $g_{tt} = -f(r)$ and $g_{rr} = f(r)^{-1}$ with $f(r) = 1 - r_s/r$.
      The parameter $r_s$ appears as the integration constant of~\eqref{eq:flux-conservation} and is identified
      operationally via its influence on admissible propagation cones.

      The derivation uses only: (i) the bounded-flux constraint (A3\textendash A4 of~\cite{Beau2026m});
      (ii) local injectivity of the projection in the exterior region; (iii) spherical symmetry and stationarity;
      (iv) absence of additional microscopic scales.
      No independent metric field equations or energy-momentum tensor are postulated.
      The detailed argument is developed in~\cite{Beau2026c}~Sections~6\textendash 8 within the Born\textendash
      Infeld dynamical picture; the present statement imports that content and reformulates it within the
      Q5b geometric setting of Section~\ref{sec:q5b}.
    \end{proof}

  \subsection{The horizon as a projection degeneracy}
    \label{sec:horizon}

    \begin{remark}[Horizon as projection degeneracy]
      \label{rem:horizon}
      At $r = r_s$, the lapse function satisfies $f(r_s) = 0$.
      In the present framework, this corresponds to a degeneracy of the principal symbol
      $\sigma_2(\mathcal{L}_{\mathrm{eff}})$ of~\eqref{eq:principal-symbol}: the operator loses hyperbolicity and the
      continuum chart ceases to be projectable.
      The horizon is therefore not a singularity of the underlying admissible structure --- it is the breakdown of
      the effective continuum description.
      This interpretation is consistent with the non-injectivity result of~\cite{Beau2026l}: at the horizon,
      distinct microscopic admissible configurations become maximally indistinguishable, saturating the projection
      entropy~$S_\Pi$.
    \end{remark}

    \begin{remark}[Relation to the Gravity paper's horizon treatment]
      \label{rem:gravity-horizon}
      The spectral entropy functional $S_\Pi[g]$ of the Gravity paper~\cite{Beau2026h} is constructed on
      Riemannian backgrounds via zeta-regularised heat-kernel methods; its extension to Lorentzian signature and to
      horizon geometries is limited by the same symbol-degeneracy mechanism identified in
      Remark~\ref{rem:horizon}.
      The two descriptions do not compete: Remark~\ref{rem:horizon} provides a structural interpretation of a
      limitation that is already implicit in the Lorentzian completion programme of~\cite{Beau2026h} and its
      thermodynamic extension~\cite{Beau2026i}.
    \end{remark}

  \section{Einstein Equations as Consistency Conditions}
  \label{sec:einstein}

  The Schwarzschild metric~\eqref{eq:schwarzschild} is the unique stationary spherically symmetric exterior solution
  of the effective geometry.
  The transition from this specific solution to the general Einstein equations is the content of the Gravity
  paper~\cite{Beau2026h}.

  \begin{remark}[Einstein equations from the spectral entropy functional]
    \label{rem:einstein}
    The Gravity paper~\cite{Beau2026h} establishes, via heat-kernel methods and zeta regularisation on Riemannian
    backgrounds, that the renormalised metric variation of the projective entropy functional
    \begin{equation}
      S_\Pi[g] \;=\; \tfrac{1}{2}\log\det{}'\mathcal{L}_\Pi
    \end{equation}
    produces, in the infrared-dominant two-derivative sector, a response proportional to the Einstein tensor
    $G_{\mu\nu}$ via the heat-kernel mechanism of Sakharov-induced gravity.
    The Einstein equations $G_{\mu\nu} = 8\pi G_N\,T_{\mu\nu}$ therefore emerge as consistency conditions of the
    emergent geometry --- not as microscopic laws of the admissible substrate.
    They express the balance between the curvature of the projection-induced metric (encoding admissibility
    constraints) and the energy-momentum content (encoding localised obstructions to the admissibility flow, such
    as the one giving~\eqref{eq:schwarzschild}).
    The companion paper on spectral thermodynamics~\cite{Beau2026i} extends this construction to a local spectral
    first law $\delta S_\Pi = \int \beta^{-1}\delta E_\Pi$.
  \end{remark}

  \begin{remark}[Chain closure]
    \label{rem:chain}
    The chain
    \begin{equation}
      \label{eq:chain}
      \Pi_q
      \;\xrightarrow{\text{Q5a}}\;
      \mathcal{L}_\Pi
      \;\xrightarrow{\text{Q5b}}\;
      g^{\mu\nu}
      \;\xrightarrow{\text{\cite{Beau2026h}}}\;
      G_{\mu\nu}
    \end{equation}
    is now closed, conditionally on the hypotheses of Section~\ref{sec:status} and on the Riemannian restriction
    of~\cite{Beau2026h}.
    The present paper supplies the middle link: the identification of the symbol-extracted object of Q5b with the
    metric appearing in the variational argument of~\cite{Beau2026h}.
  \end{remark}

  \section{Open Directions}
  \label{sec:open}

  \paragraph{Closing the hypotheses inherited from Q5a and Q5b.}
    Rigorous proofs of Hypotheses~H1,~H-w,~H-E1, Conjecture~C (Q5a) and~[H-lift] (Q5b-O1) would remove all
    conditionality from the results of the present paper.
    Q5b-O1 is rated as moderate difficulty given Q5a Theorem~T2, and is the most immediately accessible of these.

  \paragraph{Full $4 \times 4$ non-degenerate metric (Q5b-O2).}
    Theorem~\ref{thm:schwarzschild} and Remark~\ref{rem:4d} concern only the non-degenerate part of the
    sub-Riemannian symbol.
    A complete treatment requires incorporating the central $Z$-direction via sub-principal corrections
    (Beals\textendash Greiner calculus, or the tangent groupoid construction for sub-Riemannian spaces), which would
    contribute a third spatial direction to the effective geometry.
    This is Q5b-O2.

  \paragraph{Identification of $A_\tau$, $A_H$ from admissibility data (Q5b-O3).}
    The tensor entries $A_\tau$ and $A_H$ of~\eqref{eq:principal-symbol} are not yet expressed in terms of the
    Cosmochrony spectral data.
    Closing this identification would require: expressing $A_H$ in terms of the Born\textendash Infeld saturation
    constant $c_\chi$ and the spectral weight~$A$ of Q5a; and expressing $A_\tau$ in terms of the projective
    temporal ordering data of~\cite{Beau2026pto}.
    This is the first step toward an explicit metric from first principles and, through~\cite{Beau2026h}, toward a
    quantitative determination of Newton's constant~$G_N$.

  \paragraph{Cosmological implications.}
    In homogeneous regimes, the admissibility flow gives a linear growth of cumulative projected capacity $I(n)$,
    underlying the emergent Hubble law and effective cosmological expansion~\cite{Cosmochrony}.
    The precise connection between the spectral exponent $\delta_{\mathrm{pair}}$ of the O-series~\cite{Beau2026a28}
    and the observable cosmological parameters ($H_0$, $w$) is an open quantitative problem.

  \section*{Acknowledgements}

  The author acknowledges the use of large language models as a supportive tool for refining language, structure,
  and internal consistency during the development of this manuscript.
  All conceptual contributions, theoretical choices, and interpretations remain the sole responsibility of the
  author.

  \clearpage
  \bibliographystyle{cosmochrony}
  \bibliography{cosmochrony-bibliography}

\end{document}
